% Part: incompleteness
% Chapter: theories-computability
% Section: q-is-ce

\documentclass[../../../include/open-logic-section]{subfiles}

\begin{document}

\olfileid{inc}{tcp}{qce}

\olsection{$\Th{Q}$ একটি সম্পূর্ণ \printtoken{S}{c.e.} সেট}


\begin{thm}
$\Th{Q}$ !!{c.e.}, কিন্তু নির্ণেয় নয়। বস্তুত, এটি একটি সম্পূর্ণ
!!{c.e.} সেট।
\end{thm}

\begin{proof}
$\Th{Q}$ যে !!{c.e.}, তা দেখা কঠিন নয়; কারণ এটি সেই সব বাক্যের
(কোডের) সেট~$y$, যাদের $\Th{Q}$-তে $y$-এর কোনো প্রমাণ~$x$ আছে:
\[
Q = \Setabs{y}{\lexists[x][\Prf[\Th{Q}](x,y)]}.
\]
কিন্তু আমরা জানি, $\Prf[\Th{Q}](x,y)$ গণনসাধ্য (বস্তুত আদিম
পুনরাবৃত্ত), আর উপরের আকারে লেখা যায় এমন প্রতিটি সেটই c.e.

এটি একটি সম্পূর্ণ c.e.\ সেট বলার অর্থ $K \leq_m Q$ বলা, যেখানে
$K = \Setabs{x}{!A_x(x) \downarrow}$। তাই দেখাই যে $K$-কে $\Th{Q}$-তে
হ্রাস করা যায়। ক্লিনির বিধেয় $T(e,x,s)$ আদিম পুনরাবৃত্ত বলে সেটি
$\Th{Q}$-তে কোনো সূত্র~$!A_T$ দ্বারা উপস্থাপিত। তখন প্রত্যেক~$x$-এর জন্য
\begin{align*}
x \in K & \lif \lexists[s][T(x,x,s)] \\
& \lif \lexists[s][(\Th{Q} \Proves !A_T(\num x,\num x, \num s))]
  \\
& \lif \Th{Q} \Proves \lexists[s][!A_T(\num x, \num x, s)].
\end{align*}
বিপরীত দিকে, ধরি $\Th{Q} \Proves
\lexists[s][!A_T(\num x, \num x, s)]$। $\Th{Q}$-এর স্বতঃসিদ্ধগুলি
মানক মডেলে সত্য এবং প্রথম-ক্রমের নিষ্পাদন যথার্থ, তাই ওই অস্তিত্বমূলক
বাক্যটিও মানক মডেলে সত্য। অতএব কোনো স্বাভাবিক সংখ্যা~$n$-এর জন্য
$!A_T(\num x, \num x, \num n)$ সত্য। এখন $T(x,x,n)$ মিথ্যা হলে,
$!A_T$ যে $T$-কে উপস্থাপন করে তা থেকে
$\Th{Q} \Proves \lnot !A_T(\num x, \num x, \num n)$ হতো। যথার্থতা
অনুসারে এই অস্বীকারটিও মানক মডেলে সত্য হতো, যা আগের সত্যতার সঙ্গে
বিরোধ। তাই $T(x,x,n)$ সত্য, এবং সুতরাং $!A_x(x) \downarrow$।
% BN-SRC-234 TARGET-CORRECTION-BEGIN
\par\smallskip\noindent\textnormal{\small\textbf{উৎস-সংশোধন (BN-SRC-234):}
হিমায়িত ইংরেজি বিপরীত দিকটিতে $T(x,x,n)$ মিথ্যা ধরে $\Th{Q}$ যে
$\lnot !A_T(\num x,\num x,\num n)$ প্রমাণ করবে, সেই সত্য সূত্রটিকেই ভুল
করে ``মিথ্যা সূত্র'' বলেছে। বাংলা প্রমাণে মানক মডেলে $\Th{Q}$-এর যথার্থতা
স্পষ্টভাবে ব্যবহার করে $!A_T$ ও তার অস্বীকারের সত্যতার বিরোধটি রাখা হয়েছে।}
% BN-SRC-234 TARGET-CORRECTION-END

সংক্ষেপে, প্রত্যেক~$x$-এর জন্য $x$ $K$-তে থাকে যদি এবং কেবল যদি
$\Th{Q}$ সূত্র $\lexists[s][!A_T(\num x,\num x,s)]$ প্রমাণ করে। তাই
$x$-কে $\lexists[s][!A_T(\num x,\num x,s)]$ বাক্যের (একটি কোডে) পাঠায়
যে অপেক্ষক~$f$, সেটি $K$ থেকে $\Th{Q}$-তে একটি হ্রাসকরণ।
% BN-SRC-235 TARGET-CORRECTION-BEGIN
\par\smallskip\noindent\textnormal{\small\textbf{উৎস-সংশোধন (BN-SRC-235):}
হিমায়িত শেষ অনুচ্ছেদে বস্তু-ভাষার বাক্যের ভিতরে মেটা-সম্বন্ধ~$T$ বসানো
হয়েছে। বাংলা পাঠে ঠিক আগেই নির্ধারিত উপস্থাপক সূত্র~$!A_T$ রাখা হয়েছে,
যাতে হ্রাসকরণের মান সত্যিই $\Th{Q}$-এর ভাষার একটি বাক্যের কোড হয়।}
% BN-SRC-235 TARGET-CORRECTION-END
\end{proof}

\end{document}
